\documentclass[12pt,a4paper]{article}
\usepackage[utf8]{inputenc}
\usepackage[brazil]{babel}
\usepackage{amsmath,amssymb,amsfonts}
\usepackage{graphicx}
\usepackage{geometry}
\usepackage{booktabs}
\usepackage{algorithm}
\usepackage{algpseudocode}
\usepackage{hyperref}
\usepackage{xcolor}
\usepackage{listings}
\usepackage{tikz}
\usetikzlibrary{shapes,arrows,positioning}

\geometry{margin=2.5cm}

\title{\textbf{Algoritmo de Alocacao via Otimizador PL v2.0}\\
\large Sistema SIVIRA - Otimizacao de Producao}
\author{Documentacao Tecnica}
\date{Dezembro 2025}

\begin{document}

\maketitle

\tableofcontents
\newpage

\section{Introducao}

O Otimizador PL v2.0 e um sistema de otimizacao de producao que utiliza \textbf{Programacao Linear Inteira Mista (MILP)} para determinar a ordem otima de execucao de pedidos em uma linha de producao.

O sistema foi desenvolvido utilizando o solver \textbf{OR-Tools CP-SAT} do Google, que oferece alta performance para problemas de satisfacao de restricoes e otimizacao combinatoria.

\subsection{Objetivos do Sistema}

\begin{itemize}
    \item Maximizar o numero de pedidos atendidos dentro dos prazos (deadlines)
    \item Respeitar restricoes de equipamentos (capacidade finita)
    \item Garantir sequenciamento correto de atividades (gaps temporais)
    \item Otimizar a ordem de execucao considerando urgencia dos pedidos
\end{itemize}

\section{Arquitetura do Sistema}

O sistema e composto por 4 fases principais:

\begin{center}
\begin{tikzpicture}[node distance=1.2cm, auto,
    fase/.style={draw, rectangle, rounded corners, minimum width=5cm, minimum height=0.9cm, font=\small}]

    \node[fase, fill=blue!20] (f1) {Fase 1: Deteccao de Modo};
    \node[fase, fill=green!20, below=of f1] (f2) {Fase 2: Calculo Backward};
    \node[fase, fill=orange!20, below=of f2] (f3) {Fase 3: Otimizacao PL};
    \node[fase, fill=red!20, below=of f3] (f4) {Fase 4: Execucao};

    \draw[->, thick] (f1) -- (f2);
    \draw[->, thick] (f2) -- (f3);
    \draw[->, thick] (f3) -- (f4);
\end{tikzpicture}
\end{center}

\subsection{Componentes}

\begin{table}[h]
\centering
\begin{tabular}{ll}
\toprule
\textbf{Modulo} & \textbf{Funcao} \\
\midrule
\texttt{detector\_modo.py} & Detecta modo de otimizacao \\
\texttt{calculador\_horarios\_deterministicos.py} & Calculo backward scheduling \\
\texttt{modelo\_pl\_ordenacao.py} & Modelo PL com OR-Tools CP-SAT \\
\texttt{aplicador\_ordenacao.py} & Executa pedidos na ordem otimizada \\
\texttt{executor\_unificado.py} & Orquestrador principal \\
\texttt{executor\_v2.py} & Adaptador de compatibilidade \\
\bottomrule
\end{tabular}
\caption{Modulos do sistema PL v2.0}
\end{table}

\section{Formulacao Matematica}

\subsection{Variaveis de Decisao}

O modelo utiliza variaveis binarias para representar a selecao de pedidos e suas posicoes:

\begin{equation}
x_{p,j} \in \{0, 1\}
\end{equation}

Onde:
\begin{itemize}
    \item $p$ = indice do pedido ($p \in P$)
    \item $j$ = indice da janela temporal candidata ($j \in J_p$)
    \item $x_{p,j} = 1$ significa que o pedido $p$ sera executado na janela $j$
\end{itemize}

Para o modelo de ordenacao, tambem sao utilizadas:

\begin{equation}
pos_i \in \{0, 1, \ldots, n-1\}
\end{equation}

Onde $pos_i$ representa a posicao do pedido $i$ na ordem de execucao.

\subsection{Funcao Objetivo}

O objetivo e maximizar o numero total de pedidos atendidos:

\begin{equation}
\max \sum_{p \in P} \sum_{j \in J_p} x_{p,j}
\end{equation}

\subsection{Restricoes}

\subsubsection{Restricao 1: Unicidade por Pedido}

Cada pedido utiliza no maximo uma janela temporal:

\begin{equation}
\sum_{j \in J_p} x_{p,j} \leq 1, \quad \forall p \in P
\end{equation}

\subsubsection{Restricao 2: Permutacao Unica (Ordenacao)}

Para o modelo de ordenacao, todas as posicoes devem ser diferentes:

\begin{equation}
pos_i \neq pos_j, \quad \forall i \neq j
\end{equation}

Esta restricao e implementada usando a constraint \texttt{AddAllDifferent} do OR-Tools.

\subsubsection{Restricao 3: Tempo Maximo de Espera (Gap Zero)}

Para atividades consecutivas com tempo maximo de espera igual a zero:

\begin{equation}
t_{fim}^{(a)} = t_{inicio}^{(a+1)}
\end{equation}

Isso garante que as atividades sejam contiguas, sem intervalo entre elas.

\subsubsection{Restricao 4: Conflitos Temporais (Equipamentos)}

Se duas janelas se sobrepoem no tempo, nao podem ser selecionadas simultaneamente:

\begin{equation}
x_{p_1,j_1} + x_{p_2,j_2} \leq 1
\end{equation}

Para todo par $(p_1,j_1), (p_2,j_2)$ onde:
\begin{equation}
[t_1^{ini}, t_1^{fim}] \cap [t_2^{ini}, t_2^{fim}] \neq \emptyset
\end{equation}

\subsubsection{Restricao 5: Fim Obrigatorio (Deadline Rigido)}

Para pedidos com deadline obrigatorio:

\begin{equation}
x_{p,j} = 0, \quad \forall j : |t_j^{fim} - d_p| > \epsilon
\end{equation}

Onde $d_p$ e o deadline do pedido $p$ e $\epsilon = 1$ minuto (tolerancia).

\section{Backward Scheduling}

Para o cenario deterministico ($\tau_{max} = 0$), o sistema utiliza \textbf{backward scheduling} para calcular horarios fixos:

\subsection{Algoritmo}

\begin{algorithm}
\caption{Backward Scheduling}
\begin{algorithmic}[1]
\State $t_{atual} \gets deadline$
\For{cada atividade $a$ em ordem reversa}
    \State $t_{fim}^{(a)} \gets t_{atual}$
    \State $t_{inicio}^{(a)} \gets t_{fim}^{(a)} - duracao_a$
    \State $t_{atual} \gets t_{inicio}^{(a)}$
\EndFor
\end{algorithmic}
\end{algorithm}

\subsection{Representacao Visual}

\begin{center}
\begin{tikzpicture}[node distance=0.5cm]
    \node[draw, rectangle, fill=gray!30, minimum width=2cm, minimum height=0.8cm] (a1) {Ativ 1};
    \node[draw, rectangle, fill=gray!30, minimum width=2cm, minimum height=0.8cm, right=of a1] (a2) {Ativ 2};
    \node[draw, rectangle, fill=gray!30, minimum width=2cm, minimum height=0.8cm, right=of a2] (a3) {...};
    \node[draw, rectangle, fill=blue!30, minimum width=2cm, minimum height=0.8cm, right=of a3] (an) {Ativ n};
    \node[right=0.3cm of an] (dl) {\textbf{Deadline}};

    \draw[->, thick] (an) -- (a3);
    \draw[->, thick] (a3) -- (a2);
    \draw[->, thick] (a2) -- (a1);

    \node[below=0.5cm of a1] {Inicio};
    \node[below=0.5cm of an] {Fim};
\end{tikzpicture}
\end{center}

A ultima atividade termina exatamente no deadline, e as anteriores sao calculadas retroativamente.

\section{Complexidade Computacional}

\begin{table}[h]
\centering
\begin{tabular}{ll}
\toprule
\textbf{Componente} & \textbf{Complexidade} \\
\midrule
Variaveis de decisao & $O(n \cdot k)$ \\
Restricoes de unicidade & $O(n)$ \\
Restricoes de conflito & $O(n^2 \cdot k^2)$ (pior caso) \\
\bottomrule
\end{tabular}
\caption{Complexidade por componente ($n$ = pedidos, $k$ = janelas/pedido)}
\end{table}

O solver OR-Tools CP-SAT utiliza tecnicas avancadas de propagacao de restricoes e busca para encontrar solucoes otimas ou viaveis dentro do timeout configurado (padrao: 30 segundos para ordenacao, 600 segundos para modelo completo).

\section{Comparacao: Sistema Sequencial vs PL v2.0}

\begin{table}[h]
\centering
\begin{tabular}{lcc}
\toprule
\textbf{Aspecto} & \textbf{Sequencial} & \textbf{PL v2.0} \\
\midrule
Ordem de execucao & Fixa (1, 2, 3...) & Otimizada por deadline \\
Conflitos de equipamentos & Ignora & Modela explicitamente \\
Gaps temporais & Nao verifica & Restricao $\tau_{max}=0$ \\
Funcao objetivo & Nenhuma & Maximizar pedidos \\
Metodo de solucao & Greedy & MILP (CP-SAT) \\
\bottomrule
\end{tabular}
\caption{Comparacao entre metodos}
\end{table}

\section{Implementacao}

\subsection{Exemplo de Uso}

\begin{lstlisting}[language=Python, basicstyle=\ttfamily\small]
from otimizador_v2 import ExecutorV2

# Criar executor
executor = ExecutorV2()
executor.inicializar()

# Otimizar pedidos
solucao = executor.otimizar_pedidos(pedidos)

# Verificar resultados
print(f"Pedidos atendidos: {solucao.pedidos_atendidos}")
print(f"Taxa de sucesso: {solucao.taxa_sucesso}%")
\end{lstlisting}

\subsection{Parametros de Configuracao}

\begin{itemize}
    \item \textbf{timeout\_segundos}: Tempo maximo para resolucao (padrao: 600s)
    \item \textbf{inicio\_jornada}: Horario de inicio da jornada de trabalho
\end{itemize}

\section{Conclusao}

O Otimizador PL v2.0 oferece uma abordagem matematicamente fundamentada para o problema de scheduling de producao, utilizando tecnicas de otimizacao combinatoria para maximizar o atendimento de pedidos respeitando restricoes de recursos e prazos.

A arquitetura modular permite facil extensao para novos modos de operacao (como o modo FLEXIVEL planejado para fases futuras) e integracao com o sistema de menu existente.

\end{document}
